\skripsifrontchapter{PRAKATA}

Prakata pada dokumen ini merupakan contoh petunjuk penulisan. Pengguna dapat menjelaskan konteks penyusunan skripsi, menyampaikan terima kasih atas bantuan yang benar-benar diterima, serta menuliskan harapan yang sesuai dengan isi karya. Uraian tidak perlu mengulang abstrak, latar belakang, atau simpulan penelitian.

Pihak yang disebut dalam ucapan terima kasih dipilih berdasarkan kontribusinya. Nama, gelar, jabatan, dan bentuk bantuan diperiksa agar penulis tidak memberikan keterangan yang keliru. Bantuan akademik, dukungan administratif, serta dukungan pribadi dapat dijelaskan secara singkat tanpa mengungkap informasi yang tidak diperlukan dalam naskah publik.

Sebelum naskah digunakan, ganti seluruh paragraf petunjuk ini dengan prakata penulis. Tempat dan tanggal pada bagian akhir diisi sesuai penyelesaian atau ketentuan penyerahan yang berlaku. Ruang di bawah sengaja tidak memuat tanggal ataupun identitas pribadi nyata.

\vspace{1.5cm}
\begin{flushright}
  Tempat, tanggal: \rule{4.0cm}{0.4pt}\par
  \vspace{1.6cm}
  \namamahasiswa
\end{flushright}
